\section{Parité}
Les tests que nous avons faits jusqu'à présent dans les instructions
``if'' concernaient l'orientation du \robot.  En \pascal\ il est
possible de faire des tests concernant les valeurs des variables.  Les
seules variables que nous utilisons jusque là sont les compteurs de
boucles ``for''.  Il peut être utile de faire un test sur le compteur
dans le corps d'une boucle pour exécuter des actions différentes à
chaque itération.  Prenons l'exemple du dessin de la frise en forme de
double peigne ci-dessous~:

 {\raggedright 
  \includevspacePostScript[1.5cm]{FriseDoublePeigne}{1.0}{1.0}
   }

Les dents du peigne étant dessinées de gauche à droite, les
$1^{\mathrm{\grave{e}re}}$, $3^{\mathrm{\grave{e}me}}$,
$5^{\mathrm{\grave{e}me}}$, \ldots\ sont dirigées vers le bas alors
que les $2^{\mathrm{\grave{e}me}}$, $4^{\mathrm{\grave{e}me}}$,
$6^{\mathrm{\grave{e}me}}$, \ldots\ sont dirigées vers le haut.  Dans
le programme de dessin on utilise un compteur de boucle {\idt I} qui
vaut 1 quand on dessine la $1^{\mathrm{\grave{e}re}}$ dent, 2 quand on
dessine la $2^{\mathrm{\grave{e}me}}$, \ldots Par conséquent quand la
valeur du compteur de boucle est {\em impaire} (c'est-à-dire 1, 3, 5,
\ldots) il faut dessiner une dent vers le bas tandis que si la valeur
du compteur est {\em paire} (c'est-à-dire 2, 4, 6, \ldots) il faut
dessiner une dent vers le haut.  \index{Parité}Pour tester que la
valeur du compteur {\idt I} est impaire on utilise l'expression
booléenne {\idt odd$($I$)$}, le mot anglais `odd' signifiant
précisément `impair'.


{
\parindent \parindentprog
\medskip

\carrenoir
\vspace{\neglskip}\mbox{\blc  \idt  \mcl program\idt\ FriseDoublePeigne$;$
}\par 
\vspace{\neglskip}\mbox{\blc \blk\idt \mcl uses\idt \ \ldots\ $;$
}\par 
\pagebreak[3] 
\vspace{\neglskip}\mbox{\blc  \blk\idt  \mcl const\idt \ NombreDeDents\ $=$\ 30$;$}\par 
\pagebreak[3] 
\vspace{\neglskip}\mbox{\blc  \blk\idt  \mcl var\idt \ I\ $:$\ integer$;$}\par 
\vspace{\neglskip}\mbox{\blc  \idt  \mcl begin\idt }\par 
\vspace{\neglskip}\mbox{\blc  \blk\idt  ag$;$\ pe$;$}\par 
\pagebreak[3] 
\vspace{\neglskip}\mbox{\blc  \blk\idt  \mcl for\idt \ I\ $:$$=$\ 1\ \mcl to\idt \ NombreDeDents\ \mcl do\idt }\par 
\vspace{\neglskip}\mbox{\blc  \blk\blk\idt  \mcl begin\idt }\par 
\vspace{\neglskip}\mbox{\blc  \blk\blk\blk\idt  \cmt  \{\ Trait\ horizontal\ \}\idt \ av$;$}\par 
\pagebreak[3] 
\vspace{\neglskip}\mbox{\blc  \blk\blk\blk\idt  \mcl if\idt \ odd$($I$)$\ \mcl then\idt }\par 
\vspace{\neglskip}\mbox{\blc  \blk\blk\blk\blk\idt  \mcl begin\idt }\par 
\vspace{\neglskip}\mbox{\blc  \blk\blk\blk\blk\blk\idt  \cmt  \{\ Dent\ vers\ le\ bas\ \}\idt \ pqtd$;$\ av$;$\ pdt$;$\ av$;$\ pqtd$;$\ }\par 
\vspace{\neglskip}\mbox{\blc  \blk\blk\blk\blk\idt  \mcl end\idt }\par 
\pagebreak[3] 
\vspace{\neglskip}\mbox{\blc  \blk\blk\blk\idt  \mcl else\idt }\par 
\vspace{\neglskip}\mbox{\blc  \blk\blk\blk\blk\idt  \mcl begin\idt }\par 
\vspace{\neglskip}\mbox{\blc  \blk\blk\blk\blk\blk\idt  \cmt  \{\ Dent\ vers\ le\ haut\ \}\idt \ pqtg$;$\ av$;$\ pdt$;$\ av$;$\ pqtg$;$\ }\par 
\vspace{\neglskip}\mbox{\blc  \blk\blk\blk\blk\idt  \mcl end\idt $;$}\par 
\vspace{\neglskip}\mbox{\blc  \blk\blk\idt  \mcl end\idt $;$}\par 
\vspace{\neglskip}\mbox{\blc  \blk\idt  st$;$\ }\par 
\vspace{\neglskip}\mbox{\blc  \idt  \mcl end\idt $.$}\par 

\medskip }



\section{Comparaison d'entiers}

Si {\idt I} est une variable, par exemple un compteur de boucle,
l'expression booléenne {\idt $($I $>$ 10$)$} est {\idt true} (vraie)
si la valeur de {\idt I} est strictement plus grande que 10 et {\idt
false} (fausse) dans le cas contraire c'est-à-dire si la valeur de {\idt
I} est inférieure ou égale à 10.  Les symboles mathématiques $\leq$
(inférieur ou égal) et $\geq$ (supérieur ou égal) se notent $<=$ et
$>=$ en \pascal\ (car les symboles mathématiques $\leq$ et $\geq$ ne
se trouvent pas sur tous les claviers d'ordinateurs alors que $<$, $>$
et $=$ sont toujours présents).  Par exemple les expressions
booléennes $($10 $<=$ 10$)$ et $($10 $<=$ 20$)$ valent {\idt true}
(vrai) tandis que $($10 $<=$ 9$)$ vaut {\idt false} (faux).  En
\pascal\ on note l'égalité avec le symbole mathématique usuel $=$. 
Par contre le symbole mathématique $\neq$ (différent) se note $<>$. 
Par exemple $($10 $=$ 10$)$ et $($10 $<>$ 20$)$ valent {\idt true}
(vrai) tandis que $($10 $<>$ 10$)$ et $($10 $=$ 20$)$ valent {\idt
false} (faux).  On peut utiliser les expressions booléennes de
comparaison d'entiers ou de valeurs de variables entières dans le test
de l'instruction ``if''.  Par exemple la commande {\idt avf(N)} qui
fait avancer le \robot\ {\idt N} fois si la valeur de {\idt N} est
positive, le fait reculer si la valeur de {\idt N} est négative et qui
ne change pas sa position quand la valeur de {\idt N} est nulle peut
se définir comme suit~:

{
\parindent \parindentprog
\medskip

\pagebreak[3] 
\vspace{\neglskip}\mbox{\blc  \blk\idt  \mcl procedure\idt \ Avf\ $($N\ $:$\ integer$)$$;$}\par 
\pagebreak[3] 
\vspace{\neglskip}\mbox{\blc  \blk\blk\idt  \mcl var\idt \ I\ $:$\ integer$;$}\par 
\vspace{\neglskip}\mbox{\blc  \blk\idt  \mcl begin\idt }\par 
\pagebreak[3] 
\vspace{\neglskip}\mbox{\blc  \blk\blk\idt  \mcl if\idt \ $($N\ $>$\ 0$)$\ \mcl then\idt }\par 
\vspace{\neglskip}\mbox{\blc  \blk\blk\blk\idt  \mcl begin\idt }\par 
\pagebreak[3] 
\vspace{\neglskip}\mbox{\blc  \blk\blk\blk\blk\idt  \mcl for\idt \ I\ $:$$=$\ 1\ \mcl to\idt \ N\ \mcl do\idt }\par 
\vspace{\neglskip}\mbox{\blc  \blk\blk\blk\blk\blk\idt  \mcl begin\idt }\par 
\vspace{\neglskip}\mbox{\blc  \blk\blk\blk\blk\blk\blk\idt  av$;$}\par 
\vspace{\neglskip}\mbox{\blc  \blk\blk\blk\blk\blk\idt  \mcl end\idt $;$}\par 
\vspace{\neglskip}\mbox{\blc  \blk\blk\blk\idt  \mcl end\idt }\par 
\pagebreak[3] 
\vspace{\neglskip}\mbox{\blc  \blk\blk\idt  \mcl else\idt \ \mcl if\idt \ $($N\ $<$\ 0$)$\ \mcl then\idt }\par 
\vspace{\neglskip}\mbox{\blc  \blk\blk\blk\idt  \mcl begin\idt }\par 
\vspace{\neglskip}\mbox{\blc  \blk\blk\blk\blk\idt  pdt$;$\ }\par 
\pagebreak[3] 
\vspace{\neglskip}\mbox{\blc  \blk\blk\blk\blk\idt  \mcl for\idt \ I\ $:$$=$\ 1\ \mcl to\idt \ $($$-$\ N$)$\ \mcl do\idt }\par 
\vspace{\neglskip}\mbox{\blc  \blk\blk\blk\blk\blk\idt  \mcl begin\idt }\par 
\vspace{\neglskip}\mbox{\blc  \blk\blk\blk\blk\blk\blk\idt  av$;$}\par 
\vspace{\neglskip}\mbox{\blc  \blk\blk\blk\blk\blk\idt  \mcl end\idt $;$}\par 
\vspace{\neglskip}\mbox{\blc  \blk\blk\blk\blk\idt  pdt$;$\ }\par 
\vspace{\neglskip}\mbox{\blc  \blk\blk\blk\idt  \mcl end\idt $;$}\par 
\vspace{\neglskip}\mbox{\blc  \blk\idt  \mcl end\idt $;$}\par 
\medskip
 }

De manière plus générale, on peut utiliser les opérateurs $<$ (strictement inférieur), $<=$ (inférieur ou égal), 
$=$ (égal), $<>$ (différent), $>=$ (supérieur ou égal) et $>$ (strictement supérieur) pour comparer les valeurs
de deux expressions entières comme par exemple {\idt $($I $<=$ L$)$} dans le programme suivant qui trace le serpentin
ci-dessous~:

 \includevspacePostScript[1.5cm]{Serpentin}{0.95}{0.95}
 \input Serpentin.tex

\section{Négation, conjonction et disjonction}

La {\em négation} \index{Négation} d'une expression booléenne $E$ s'écrit $(${\mcl not} $E)$, le mot anglais {\mcl not} signifiant `non'. Si la
valeur de $E$ est {\idt true} alors celle de $(${\mcl not} $E)$ est {\idt false} tandis que si la valeur de $E$ est {\idt false} alors
celle de $(${\mcl not} $E)$ est {\idt true}. Par exemple  {\idt $(${\mcl not} odd$($5$))$} vaut {\idt false} puisque
{\idt odd$($5$)$} vaut {\idt true}, {\idt 5} étant impair. Ceci peut se résumer par la table de vérité suivante~:

\begin{center}
\begin{tabular}{|c|c|} \hline
$E$ & $(${\mcl not}$E)$\\ \hline \hline
{\idt true} & {\idt false}\\ \hline
{\idt false} & {\idt true}\\ \hline
\end{tabular}
\end{center}

La {\em conjonction} \index{Conjonction} de deux expressions booléennes $E_{1}$ et $E_{2}$ s'écrit $(E_{1}$ {\mcl and} $E_{2})$, 
le mot anglais {\mcl and} signifiant `et'. La valeur de $(E_{1}$ {\mcl and} $E_{2})$ est {\idt true} (vrai) si et
seulement si la valeur de $E_{1}$ et celle de $E_{2}$ sont {\idt true} (vrai). Si l'une des expressions
booléennes $E_{1}$ ou 
$E_{2}$ vaut {\idt false} (faux) alors la conjonction $(E_{1}${\mcl and} $E_{2})$ vaut également {\idt false} (faux).
Par exemple  $((5 >= 0)$ {\mcl and} {\idt odd$($5$))$} vaut {\idt true} puisque
{\idt 5} est positif et impair. Ceci peut se résumer par la table de vérité suivante~:

\begin{center}
\begin{tabular}{|c|c|c|} \hline
$E_{1}$ & $E_{2}$ & $(E_{1}${\mcl and} $E_{2})$\\ \hline \hline
{\idt true} & {\idt true} & {\idt true}\\ \hline
{\idt true} & {\idt false} & {\idt false}\\ \hline
{\idt false} & {\idt true} & {\idt false}\\ \hline
{\idt false} & {\idt false} & {\idt false}\\ \hline
\end{tabular}
\end{center}

La {\em disjonction} \index{Disjonction} de deux expressions booléennes $E_{1}$ et $E_{2}$ s'écrit $(E_{1}$ {\mcl or} $E_{2})$, 
le mot anglais {\mcl or} signifiant `ou'. La valeur de $(E_{1}$ {\mcl or} $E_{2})$ est {\idt true} (vrai) si et
seulement si la valeur de $E_{1}$ ou celle de $E_{2}$ est {\idt true} (vrai). Si les valeurs de $E_{1}$ et de
$E_{2}$ sont  {\idt false} (faux) alors la disjonction $(E_{1}$ {\mcl or} $E_{2})$ vaut également {\idt false} (faux).
Par exemple  $((4 >= 0)$ {\mcl or} {\idt odd$($4$))$} vaut {\idt true} puisque
{\idt 4} est positif et bien que {\idt 4} soit pair. Ceci peut se résumer par la table de vérité suivante~:

\begin{center}
\begin{tabular}{|c|c|c|} \hline
$E_{1}$ & $E_{2}$ & $(E_{1}${\mcl or} $E_{2})$\\ \hline \hline
{\idt true} & {\idt true} & {\idt true}\\ \hline
{\idt true} & {\idt false} & {\idt true}\\ \hline
{\idt false} & {\idt true} & {\idt true}\\ \hline
{\idt false} & {\idt false} & {\idt false}\\ \hline
\end{tabular}
\end{center}

La conjonction et la disjonction peuvent se réaliser simplement avec un circuit électrique comme on le voit à la figure
\ref{fig:andor} 
ci-dessous
%(page \pageref{fig:andor}) 
avec la con\-ven\-tion que {\idt true} représente l'interrupteur fermé
et l'ampoule allumée (le courant passe) et {\idt false} représente l'interrupteur ouvert
et l'ampoule éteinte (le courant ne passe pas).

\begin{figure}[htbp]
 \includevspacePostScript[4.3cm]{CircuitsAndOr}{0.95}{0.95}
 \caption{Circuits électriques réalisant la conjonction et\label{fig:andor}}
\centering{la disjonction logiques}
\end{figure}

Une disjonction est utilisée dans le programme suivant pour tester si le \robot\ est orienté au nord ou à
l'est~:

\vspace{0.3cm}
 {\raggedright 
 \includevspacePostScript[0.5cm]{FriseEnOr}{1.0}{1.0}

{
\parindent \parindentprog
\medskip

\carrenoir
\vspace{\neglskip}\mbox{\blc  \idt  \mcl program\idt\ FriseEnOr$;$
}\par 
\vspace{\neglskip}\mbox{\blc \blk\idt \mcl uses\idt \ \ldots\ $;$
}\par 
\pagebreak[3] 
\vspace{\neglskip}\mbox{\blc  \blk\idt  \mcl const\idt \ Longueur\ $=$\ 31$;$}\par 
\pagebreak[3] 
\vspace{\neglskip}\mbox{\blc  \blk\idt  \mcl var\idt \ I\ $:$\ integer$;$}\par 
\vspace{\neglskip}\mbox{\blc  \idt  \mcl begin\idt }\par 
\vspace{\neglskip}\mbox{\blc  \blk\idt  ag$;$}\par 
\pagebreak[3] 
\vspace{\neglskip}\mbox{\blc  \blk\idt  \mcl for\idt \ I\ $:$$=$\ 1\ \mcl to\idt \ Longueur\ \mcl do\idt }\par 
\vspace{\neglskip}\mbox{\blc  \blk\blk\idt  \mcl begin\idt \ }\par 
\pagebreak[3] 
\vspace{\neglskip}\mbox{\blc  \blk\blk\blk\idt  \mcl if\idt \ $($tn\ \mcl or\idt \ te$)$\ \mcl then\idt }\par 
\vspace{\neglskip}\mbox{\blc  \blk\blk\blk\blk\idt  \mcl begin\idt }\par 
\vspace{\neglskip}\mbox{\blc  \blk\blk\blk\blk\blk\idt  vd$;$\ }\par 
\vspace{\neglskip}\mbox{\blc  \blk\blk\blk\blk\idt  \mcl end\idt }\par 
\pagebreak[3] 
\vspace{\neglskip}\mbox{\blc  \blk\blk\blk\idt  \mcl else\idt }\par 
\vspace{\neglskip}\mbox{\blc  \blk\blk\blk\blk\idt  \mcl begin\idt }\par 
\vspace{\neglskip}\mbox{\blc  \blk\blk\blk\blk\blk\idt  avf$($$-$1$)$$;$\ vg$;$\ pqtg$;$\ }\par 
\vspace{\neglskip}\mbox{\blc  \blk\blk\blk\blk\idt  \mcl end\idt $;$}\par 
\vspace{\neglskip}\mbox{\blc  \blk\blk\idt  \mcl end\idt $;$}\par 
\vspace{\neglskip}\mbox{\blc  \blk\idt  st$;$}\par 
\vspace{\neglskip}\mbox{\blc  \idt  \mcl end\idt $.$}\par 

\medskip }


  }

\section{Expressions booléennes}

La table ci-dessous résume les règles d'écriture et les règles de calcul de la valeur
des expressions booléennes en \pascal~:

\index{Expression booléenne} \index{Test} 

\begin{center}
\protect\frenchspacing
\begin{tabular}{@{}|l|p{6.8cm}|@{}} \hline
\multicolumn{2}{|l|}{\em Expressions booléennes~:} \\ \hline \hline
{\idt true} & Valeur `vrai'.\\
{\idt false} & Valeur `faux'.\\ \hline
{\idt odd}$(${\small\rm Expr. entière}$)$ & {\idt true} si et seulement si la valeur de l'{\small\rm Expr. entière} est impaire.\\ \hline
$(${\small\rm Expr. entière$_{1}$} $<=$ {\small\rm Expr. entière$_{2}$}$)$& {\idt true} si et seulement si la valeur de l'{\small\rm Expr. entière$_{1}$} est 
                                 inférieure ou égale à celle de l'{\small\rm Expr. entière$_{2}$}.\\
$(${\small\rm Expr. entière$_{1}$} $<$ {\small\rm Expr. entière$_{2}$}$)$& {\idt true} si et seu\-le\-ment si la va\-leur de l'{\small\rm Expr. entière$_{1}$} est 
                                 strictement inférieure à celle de l'{\small\rm Expr. entière$_{2}$}.\\
$(${\small\rm Expr. entière$_{1}$} $=$ {\small\rm Expr. entière$_{2}$}$)$& {\idt true} si et seu\-le\-ment si la valeur de l'{\small\rm Expr. entière$_{1}$} est 
                                 égale à celle de l'{\small\rm Expr. entière$_{2}$}.\\
$(${\small\rm Expr. entière$_{1}$} $<>$ {\small\rm Expr. entière$_{2}$}$)$& {\idt true} si et seu\-le\-ment si la valeur de l'{\small\rm Expr. entière$_{1}$} est 
                                 différente de celle de l'{\small\rm Expr. entière$_{2}$}.\\ \hline
\end{tabular}

\begin{tabular}{@{}|l|p{6.8cm}|@{}} \hline
$(${\small\rm Expr. entière$_{1}$} $>$ {\small\rm Expr. entière$_{2}$}$)$& {\idt true} si et seu\-le\-ment si la va\-leur de l'{\small\rm Expr. entière$_{1}$} est 
                                 strictement supérieure à celle de l'{\small\rm Expr. entière$_{2}$}.\\
$(${\small\rm Expr. entière$_{1}$} $>=$ {\small\rm Expr. entière$_{2}$}$)$& {\idt true} si et seulement si la valeur de l'{\small\rm Expr. entière$_{1}$} est 
                                supérieure ou égale à celle de l'{\small\rm Expr. entière$_{2}$}.\\ \hline
$(${\idt not} {\small\rm Expr. bool.}$)$ & {\idt true} si et seulement si la valeur de l'{\small\rm Expr. bool.} est 
                                égale à {\idt false}.\\ \hline
$(${\small\rm Expr. bool.$_{1}$} {\idt and} {\small\rm Expr. bool.$_{2}$}$)$& {\idt true} si et seulement si les valeurs de l'{\small\rm Expr. bool.$_{1}$} et 
                                 de l'{\small\rm Expr. bool.$_{2}$} sont égales à {\idt true}.\\
$(${\small\rm Expr. bool.$_{1}$} {\idt or} {\small\rm Expr. bool.$_{2}$}$)$& {\idt true} si et seulement si la valeur de l'{\small\rm Expr. bool.$_{1}$} et 
                                 ou celle de l'{\small\rm Expr. bool.$_{2}$} est égale à {\idt true}.\\ \hline
\end{tabular}
\end{center}

\noindent \rule{137mm}{0.1mm} %---------------------------------------------
\begin{exercices}{Écrire des programmes \pascal\  utilisant des tests de parité pour réaliser 
les dessins suivants~:

\tiret  la serpentine pyramidale et la grille de la figure \ref{fig:serpyrgricr} ci-dessous$\:$;
%(page \pageref{fig:serpyrgricr})$\:$;

\tiret   le pavage octogonal et le pavage d'hexagones entrelacés de la figure \ref{fig:pavocthex} (page \pageref{fig:pavocthex}).

\noindent $\Box$}\end{exercices}
\begin{figure}[htbp]
  \includevspacePostScript[6.0cm]{SerpPyr+GrilleDeCr}{0.95}{0.95}
 \caption{Serpentine pyramidale et grille de croix}\label{fig:serpyrgricr}
\end{figure}
\begin{figure}[htbp]
\includevspacePostScript[5.4cm]{PavageOctoHexa}{0.95}{0.95}
 \caption{Pavages d'octogones juxtaposés et d'hexagones entrelacés}\label{fig:pavocthex}
\end{figure}
\begin{exercices}{Écrire des programmes \pascal\  utilisant des expressions boolé\-en\-nes de comparaison de
valeurs entières pour réaliser les dessins suivants (le deuxième est difficile)~:

\tiret  les chiffres digitaux de la figure \ref{fig:alphabet} (page \pageref{fig:alphabet}) où le chiffre à dessiner
est passé en paramètre de la procédure de dessin$\:$;

\tiret  le pavage de roues dentées de la figure \ref{fig:pavrd} ci-dessous$\:$;

\tiret  le tapis égyptien de la figure \ref{fig:tapegyp} (page \pageref{fig:tapegyp}).

\noindent $\Box$}\end{exercices}
\begin{figure}[htbp]
 \includevspacePostScript[5.7cm]{PavageRoueDentee}{0.95}{0.95}
 \caption{Pavage de roues dentées}\label{fig:pavrd}
\end{figure}
\begin{figure}[htbp]
 \includevspacePostScript[7.2cm]{TapisEgyptien}{0.95}{0.95}
 \caption{Tapis égyptien\label{fig:tapegyp}}
\end{figure}
\begin{exercices}{Écrire des programmes \pascal\  utilisant des négations, conjonctions ou disjonctions d'expressions 
booléen\-nes pour réaliser les dessins suivants (le deuxième est difficile)~:

\tiret  le carré zébré de la figure \ref{fig:carzeb} (page \pageref{fig:carzeb})$\:$;

\tiret  la mosaïque de l'Alhambra représentée à la figure \ref{fig:mosaic6} (page \pageref{fig:mosaic6}).

\noindent $\Box$}\end{exercices}
\begin{figure}[htbp]
 \includevspacePostScript[5.8cm]{CarreZebree}{0.95}{0.95}
 \caption{Carré zébré}\label{fig:carzeb}
\end{figure}
\begin{figure}[htbp]
 \includevspacePostScript[5.7cm]{MosaiqueAlhambra6}{0.95}{0.95}
 \caption{Mosaïque de l'Alhambra}\label{fig:mosaic6}
\end{figure}
%%%%%%%%%%%%%%%%%%%%%%%%%%%%%%%%%%%%%%%%%%%%%%%%%%%%%%%%%%
\begin{corrige}{
Nous donnons un programme, les autres se trouvant sur disquette~:
\input SerpentinePyramidale.tex
\noindent $\Box$}\end{corrige}
\begin{corrige}{
Nous donnons un programme, le premier et le trois\-ième se trouvant sur disquette~:
\input PavageRouesDentees.tex
\noindent $\Box$}\end{corrige}
\begin{corrige}{
Nous donnons un programme, le premier se trouvant sur disquette~:
\input MosaiqueAlhambra6.tex
\noindent $\Box$}\end{corrige}

%%%%%%%%%%%%%%%%%%%%%%%%%%%%%%%%%%%%%%%%%%%%%%%%%%%%%%%%%%%%%%%%%%%%%%%%%%%
\noindent \framebox[\textwidth]{
\begin{minipage}{\myparboxwidth}
\setlength{\parindent}{2em}
\small \sf
\vspace{\abovedisplayskip}
\section*{Addenda}
 
 En \pascal\ les priorités des opérateurs entiers et booléens sont définies comme suit~:
 
 \vspace{\baselineskip}
 
 \index{Priorité des opérateurs}
\begin{center}
\begin{tabular}{|l|l|l|} \hline
\multicolumn{3}{|l|}{\em Priorité des opérateurs entiers et booléens~:} \\ \hline \hline
$-$ (unaire), {\mcl not}                               & opérateurs unaires & priorité forte \\ \hline
$\ast$, {\mcl div}, {\mcl mod}, {\mcl and} & opérateurs multiplicatifs & priorité mi-forte \\ \hline
$+$ $-$ (binaires), {\mcl or}  & opérateurs additifs & priorité mi-faible \\ \hline
$<$, $<=$, $=$, $<>$, $>=$, $>$  & opérateurs de relation & priorité faible \\ \hline
\end{tabular}
\end{center}

\vspace{\baselineskip}

 \noindent Par exemple l'expression booléenne `{\idt  true\ \mcl and\idt \ false\ \mcl or\idt \ false}' s'é\-va\-lue com\-me
 `{\idt  $(($true\ \mcl and\idt \ false$)$\ \mcl or\idt \ false$)$}' et est donc égale à {\idt false}.

\vspace{\abovedisplayskip}
\end{minipage}
}
%%%%%%%%%%%%%%%%%%%%%%%%%%%%%%%%%%%%%%%%%%%%%%%%%%%%%%%%%%%%%%%%%%%%%%%%%%%% 
 
%%%%%%%%%%%%%%%%%%%%%%%%%%%%%%%%%%%%%%%%%%%%%%%%%%%%%%%%%%%%%%%%%%%%%%%%%%%
\noindent \framebox[\textwidth]{
\begin{minipage}{\myparboxwidth}
\setlength{\parindent}{2em}
\small \sf
\vspace{\abovedisplayskip}
\indent  Ce choix des priorités des opérateurs a été fait pour éviter un trop grand nombre de niveaux. Cependant
 cette convention conduit à certaines écritures peu naturelles. Par exemple `{\idt  0\ $<=$\ I\ \mcl and\idt \ I\ $<$\ 100}'
s'évalue comme `{\idt  $(($0\ $<=$\ $($I\ \mcl and\idt \ I$))$\ $<$\ 100$)$}' qui est incorrecte si {\idt I} est une
variable entière puisque la conjonction n'est définie en \pascal\ que pour les booléens. Dans certains langages de
programmation cette expression `{\idt  0\ $<=$\ I\ \mcl and\idt \ I\ $<$\ 100}' s'évalue
comme `{\idt  $(($0\ $<=$\ I$)$\ \mcl and\idt \ $($I\ $<$\ 100$))$}'
 qui est correcte et plus proche des conventions mathématiques usuelles. Pour éviter les confusions et les erreurs nous 
employons systématiquement des parenthèses.

Le terme ``booléen'' ou ``booléien'' ({\idt boolean} en anglais) vient de Boole\footnote{\ \href{https://fr.wikipedia.org/wiki/George_Boole}{George Boole}, mathématicien britannique, 1815--1864.},
 créateur de la logique mathématique moderne.
 
\vspace{\abovedisplayskip}
\end{minipage}
}
%%%%%%%%%%%%%%%%%%%%%%%%%%%%%%%%%%%%%%%%%%%%%%%%%%%%%%%%%%%%%%%%%%%%%%%%%%%% 
\newpage                 % Si on termine le chapitre par une page blanche
\thispagestyle{empty}
